%%%%%%%%%%%%%%%%%%%%%%%%%%%%%%%%%%%%%%%%%
% Beamer Presentation
% LaTeX Template
% Version 1.0 (10/11/12)
%
% This template has been downloaded from:
% http://www.LaTeXTemplates.com
%
% License:
% CC BY-NC-SA 3.0 (http://creativecommons.org/licenses/by-nc-sa/3.0/)
%
%%%%%%%%%%%%%%%%%%%%%%%%%%%%%%%%%%%%%%%%%

%-------------------------------------------------------------------------------
%	PACKAGES AND THEMES
%-------------------------------------------------------------------------------

\documentclass{beamer}
\usepackage{xcolor}
\usepackage{graphicx}
\usepackage{tikz}
\usepackage{tcolorbox}
\usepackage{listings}
\usepackage{multicol}

\DeclareMathOperator{\diag}{diag}

\definecolor{applegreen}{rgb}{0.55, 0.71, 0.0}
\definecolor{blue(ncs)}{rgb}{0.0, 0.45, 0.60}
\definecolor{burgundy}{rgb}{0.5, 0.0, 0.13}

\definecolor{cadet}{rgb}{0.33, 0.41, 0.47}
\definecolor{airforceblue}{rgb}{0.36, 0.54, 0.66}

\mode<presentation> {

\usetheme{CambridgeUS}

\usecolortheme{wolverine}

\definecolor{gold}{HTML}{D4A017}
\definecolor{darkgold}{HTML}{B7950B}

\setbeamercolor{palette primary}{bg=cadet,fg=white}
\setbeamercolor{palette secondary}{bg=airforceblue,fg=white}
\setbeamercolor{palette tertiary}{bg=black,fg=white}
\setbeamercolor{palette quaternary}{bg=cadet,fg=white}

\setbeamercolor{frametitle}{bg=airforceblue,fg=white}

\setbeamercolor{section number projected}{bg=black,fg=cadet}
\setbeamercolor{item}{fg=black,bg=cadet}

\setbeamertemplate{page number in head/foot}[framenumber]

\lstset{basicstyle=\ttfamily\footnotesize,breaklines=true}
}

\usepackage{graphicx} % Allows including images
\usepackage{booktabs} % Allows the use of \toprule, \midrule and \bottomrule in tables

%----------------------------------------------------------------------------------------
%	TITLE PAGE
%----------------------------------------------------------------------------------------

\title[AI for RSEng]{AI for Research Software Engineering
} % The short title appears at the bottom of every slide, the full title is only on the title page

\author{Jeremy L Thompson} % Your name
\institute[Helmholz Hereon] % Your institution as it will appear on the bottom of every slide, may be shorthand to save space
{Helmholz-Zentrum Hereon \\ % Your institution for the title page
\medskip
\textit{jeremy@jeremylt.org} % Your email address
}
\date{14 July 2026} % Date, can be changed to a custom date

\begin{document}

\begin{frame}
\titlepage % Print the title page as the first slide
\end{frame}

%-------------------------------------------------------------------------------
\section{Overview}
%-------------------------------------------------------------------------------

\begin{frame}
\begin{center}
\frametitle{Overview}

Research Software Engineers must respect both\\\textbf{software development best practices} and \textbf{research ethics}

~\\

\textbf{Not all development best practices} for LLM (AI) based tooling\\\textbf{map easily} from industry to research software development

~\\

\textbf{Research ethics impose additional constraints} beyond\\industry software development considerations

~\\

\tcbox[colback=white,colframe=purple,size=fbox,boxrule=0.8pt,on line]{Key takeaway: \textbf{LLMs are not a magic bullet}}

\end{center}
\end{frame}

%------------------------------------------------

\begin{frame}
\begin{center}
\frametitle{Software as First Class Artifacts}

Many academic groups (San Francisco DORA, JOSS, etc) advocate for\\research software as first class artifacts (like papers)\\

~\\


\includegraphics[height=2cm]{sfdora.png} \hspace{10mm} 
\includegraphics[height=2cm]{joss.png}\\
\tcbox[colback=white,colframe=red,size=fbox,boxrule=0.8pt,on line]{\href{https://sfdora.org}{https://sfdora.org}} \hspace{25mm} \tcbox[colback=white,colframe=red,size=fbox,boxrule=0.8pt,on line]{\href{https://joss.theoj.org}{https://joss.theoj.org}}

~\\

RSEs need to ensure \textbf{quality}, \textbf{clarity}, \textbf{testing},\\and \textbf{documentation} to deliver first class software

\end{center}
\end{frame}

%------------------------------------------------

\begin{frame}
\frametitle{Overview} % Table of contents slide, comment this block out to remove it
\tableofcontents % Throughout your presentation, if you choose to use \section{} and \subsection{} commands, these will automatically be printed on this slide as an overview of your presentation
\end{frame}

%-------------------------------------------------------------------------------
\section{Industry Best Practices}
%-------------------------------------------------------------------------------

\begin{frame}
\begin{center}
\frametitle{Industry Best Practices}

LLM based coding tools are subject to GIGO

~\\


\includegraphics[height=2cm]{garbage.png} \hspace{10mm} 
\includegraphics[height=2cm]{garbage.png}\\
Garbage In \hspace{3mm} $\rightarrow$ \hspace{3mm} Garbage Out

~\\

Best practices for input (codebases, planning) and output (process, review)

~\\

'Healthier' codebases more likely to result in higher quality LLM output

~\\

Best practices for human development $\rightarrow$ good codebases for LLM

~\\

High developer skill required for higher quality LLM output

\end{center}
\end{frame}

%------------------------------------------------

\begin{frame}
\begin{center}
\frametitle{Caveats}

Quality research on LLM best practices hard to identify

~\\

\begin{itemize}

\item Massive amounts of money in LLM research

~\\

\item Profit motives can complicate research context

~\\

\item Modern research more likely to include positive, definitive wording
(LLM are prone to this as well)

~\\

\item Included impressions from professional discussions, blogs, etc

\end{itemize}

\end{center}
\end{frame}

%------------------------------------------------

\begin{frame}
\begin{center}
\frametitle{Code Health}

CodeHealth™ is a quality metric used by CodeScene (paid software suite)

~\\

Corresponds to many 'code smell' concerns

~\\

\begin{itemize}

\item Code cohesion, complexity

~\\

\item Adherence to DRY

~\\

\item Object and method complexity

\end{itemize}

~\\

Higher CodeHealth $\rightarrow$ fewer failing tests after LLM refactoring

\end{center}
\end{frame}

%------------------------------------------------

\begin{frame}
\begin{center}
\frametitle{Personal Discussions}

Personal discussions and blogs largely align with CodeScene research

~\\


\includegraphics[height=1.5cm]{discussion.png}

~\\

Skilled developers needed to monitor/fix LLM output quality

~\\

A clear plan and 'definition of done' before development helps greatly

~\\

Overall design and style consistency of particular note

~\\

High developer skill required for higher quality LLM output

\end{center}
\end{frame}

%------------------------------------------------

\begin{frame}
\begin{center}
\frametitle{High Skill Needed}

RSE organizations frequently list skills and training as a weakness

~\\

\tcbox[colback=white,colframe=purple,size=fbox,boxrule=0.8pt,on line]{\includegraphics[height=3.5cm]{derse-concerns.png}}
deRSE list of known community issues

~\\

LLM is not a replacement for professional software skills

~\\

Unfortunately, non-careful LLM usage can also inhibit skill formation

\end{center}
\end{frame}

%------------------------------------------------

\begin{frame}
\begin{center}
\frametitle{De-skilling}

LLM present risk of de-skilling when too much cognition is offloaded

~\\

\includegraphics[height=3.5cm]{deskilling.png}\\
Skill erosion loop (Rinta-Kahila et al.)

~\\


Research software as first class artifacts requires clarity, correctness

~\\

De-skilling presents an important risk to clarity and correctness for RSE

\end{center}
\end{frame}

%------------------------------------------------

\begin{frame}
\begin{center}
\frametitle{Summary}

RSEs need \textbf{training in development best practices} to best use LLM

~\\

Skills and training is a \textbf{known weakness} for RSEs

~\\

LLM usage can result in \textbf{de-skilling} or \textbf{inhibit skill formation}

~\\

\textbf{Curating research software like papers} parallels industry best practices

~\\
~\\

\tcbox[colback=white,colframe=purple,size=fbox,boxrule=0.8pt,on line]{Recommendation: Ensure good development skills before using LLM}

\end{center}
\end{frame}

%-------------------------------------------------------------------------------
\section{Code Review}
%-------------------------------------------------------------------------------

\begin{frame}
\begin{center}
\frametitle{Code Review}

LLM based code review is a promising application

~\\

\includegraphics[height=4cm]{copilot.png}\\

~\\

Code review is high skill activity, so de-skilling risk still applies

~\\

Note: LLM based code review appears to consume many tokens

\end{center}
\end{frame}

%------------------------------------------------

\begin{frame}
\begin{center}
\frametitle{Options}

Some popular LLM based code review tools:

~\\

\begin{itemize}

\item GitHub Copilot

~\\

\item Coderabbit

~\\

\item Cursor

~\\

\item Claude Code

~\\

\item And many, many more

\end{itemize}

\end{center}
\end{frame}

%------------------------------------------------

\begin{frame}
\begin{center}
\frametitle{Caveats}

Total token usage and costs per token continue to rise

~\\

Most popular models controlled by large US tech companies

~\\


\includegraphics[height=.75cm]{usa.png} \hspace{10mm} 
\includegraphics[height=.75cm]{claudecode.png} \hspace{5mm} 
\includegraphics[height=.75cm]{openai.png} \hspace{5mm} 
\includegraphics[height=.75cm]{cursor.png} \hspace{5mm} 
\includegraphics[height=.75cm]{gemini.png} \hspace{5mm} 
\includegraphics[height=.75cm]{githubcopilot.png}\\

~\\

SLMs (Small Language Models) on local hardware possible alternative

~\\

SLM options: Qwen, Gemma, Phi, SmolLM, Ministral, etc

~\\

Offers data control, no API costs, but less context and complexity

\end{center}
\end{frame}

%-------------------------------------------------------------------------------
\section{Research Ethics}
%-------------------------------------------------------------------------------

\begin{frame}
\begin{center}
\frametitle{Relevant Points}

Three key ethics points:

~\\

\begin{itemize}

\item Research requires \textbf{transparency} in methods, including tools

~\\

\item Software must remain \textbf{legitimate and citable} research

~\\

\item Must \textbf{not plagiarize} work of others, even inadvertently

\end{itemize}

~\\

~\\

\tcbox[colback=white,colframe=purple,size=fbox,boxrule=0.8pt,on line]{See European Code of Conduct for Research Integrity, among others}

\end{center}
\end{frame}

%------------------------------------------------

\begin{frame}
\begin{center}
\frametitle{Disclosure}

Disclosure is a natural requirement

~\\

\begin{itemize}

\item Disclosure is at the heart of science

~\\

\item Disclosure improves reproducibility

~\\

\item Successful methods can be re-used

~\\

\item Unsuccessful methods can be avoided

\end{itemize}

~\\

BLUF: If you are uncomfortable stating what you did\\
the software/research is not ready to release and share

\end{center}
\end{frame}

%------------------------------------------------

\begin{frame}
\begin{center}
\frametitle{Open Source Policies}

Disclosure aligns with Open Source policies permitting LLM

~\\

\begin{itemize}

\item Many open source project have LLM policies

~\\

\item Disclosure usually required if LLM output permitted

~\\

\item Disclosure can improve code review

\end{itemize}

~\\
~\\

See also Melissa Weber Mendonça's list of open source policies:\\
\tcbox[colback=white,colframe=red,size=fbox,boxrule=0.8pt,on line]{\href{https://github.com/melissawm/open-source-ai-contribution-policies}{https://github.com/melissawm/open-source-ai-contribution-policies}}

\end{center}
\end{frame}

%------------------------------------------------

\begin{frame}
\begin{center}
\frametitle{Plagiarism}

Dr Tim Davis prompted Copilot to generate his copyrighted CSparse code\\
without attribution and violating LGPL licensing

~\\

\includegraphics[height=4.2cm]{TimDavisTweet.png}

~\\

Researchers are ultimately responsible for the code they create

~\\

Accidental plagiarism via LLM is still an ethics violation

\end{center}
\end{frame}

%------------------------------------------------

\begin{frame}
\begin{center}
\frametitle{Citation Errors}

GPT-4 has fabricated or generated erroneous citations

~\\

Generated citations can have slightly flawed data

~\\

LLM output can omit specific training data related to the prompt

~\\
~\\

RSE should guard against research process inversion:\\research $\rightarrow$ reach conclusions {\color{green} \checkmark}\\conclusions $\rightarrow$ find research to support {\color{red} $\times$}

\end{center}
\end{frame}

%------------------------------------------------

\begin{frame}
\begin{center}
\frametitle{Citation Bias}

LLM generated citations reinforce bias in citations

\end{center}
\end{frame}

%------------------------------------------------

\begin{frame}
\begin{center}
\frametitle{Summary}

\textbf{Easy to accidentally plagiarize} with heavy usage

~\\

LLM generated citations can be \textbf{fabricated, flawed, or reinforce bias}

~\\
~\\

\tcbox[colback=white,colframe=purple,size=fbox,boxrule=0.8pt,on line]{\textbf{RSEs are ultimately responsible for our work}, no matter the process}

\end{center}
\end{frame}

%-------------------------------------------------------------------------------
\section{Summary}
%-------------------------------------------------------------------------------

\begin{frame}
\begin{center}
\frametitle{Summary}

\tcbox[colback=white,colframe=red,size=fbox,boxrule=0.8pt,on line]{Key takeaway: \textbf{LLMs are not a magic bullet}}

~\\
~\\

RSEs who let LLM code for them will have bad code

~\\

\tcbox[colback=white,colframe=purple,size=fbox,boxrule=0.8pt,on line]{We need to ensure good development skills before integrating LLM}

~\\
~\\

RSEs who let LLM research for them will violate ethics

~\\

\tcbox[colback=white,colframe=purple,size=fbox,boxrule=0.8pt,on line]{We are ultimately responsible for our work, no matter the process}

\end{center}
\end{frame}

%------------------------------------------------

\begin{frame}
\begin{center}
\frametitle{Questions? Comments?}

Community comments, resources, and questions invited

\tcbox[colback=white,colframe=red,size=fbox,boxrule=0.8pt,on line]{\href{https://github.com/jeremylt/personal_website/pull/3}{https://github.com/jeremylt/personal\_website/pull/3}}

~\\
~\\

\tcbox[colback=white,colframe=purple,size=fbox,boxrule=0.8pt,on line]{Key takeaway: \textbf{LLMs are not a magic bullet}}

~\\

\tcbox[colback=white,colframe=purple,size=fbox,boxrule=0.8pt,on line]{\textbf{RSE need good development skills} before integrating LLM}

~\\

\tcbox[colback=white,colframe=purple,size=fbox,boxrule=0.8pt,on line]{\textbf{RSEs are ultimately responsible for our work}, no matter the process}

\end{center}
\end{frame}

%-------------------------------------------------------------------------------

\begin{frame}[noframenumbering]
\titlepage % Print the title page
\end{frame}

%-------------------------------------------------------------------------------

\end{document}

%-------------------------------------------------------------------------------
